%%% ВЫБОР УСТРОЙСТВА %%%

\IfStrEqCase{\devicecode}{%
{ЮТКБ.656122.611-05.02 РЭ}{%
% код для этого типа
\def\AbzazOne{
\item
Алгоритм работы ЗНФ В1 приведен на рисунке \ref{pict:ZNF}, алгоритм работы ЗНФ В2 выполнен аналогично. Уставки ЗНФ В1 приведены в таблице \ref{app_set:znf1}. Уставки ЗНФ В2 приведены в таблице \ref{app_set:znf2}.
}
}%
{ЮТКБ.656122.611-11.01 РЭ}{%ШР
% код для этого типа
\def\AbzazOne{
\item
Алгоритм работы ЗНФ приведен на рисунке \ref{pict:ZNF}. Уставки ЗНФ приведены в таблице \ref{app_set:znf}.
}
}%
{ЮТКБ.656122.611-09.01 РЭ}{%БСК
% код для этого типа
\def\AbzazOne{
\item
Алгоритм работы ЗНФ приведен на рисунке \ref{pict:ZNF}. Уставки ЗНФ приведены в таблице \ref{app_set:znf}.
}
}%
}
[
\def\AbzazOne{??}
]


\color{unimain}{\subsection{Защита от непереключения фаз}}
\color{black}

\begin{enumerate}

\item
Защита от непереключения фаз (ЗНФ) предназначена для обнаружения расхождения полюсов силового выключателя с пофазным приводом при включении или отключении.   


\item
Ввод в работу ЗНФ осуществляется с помощью программного переключателя \textbf{«ЗНФ»}. Факт расхождения полюсов может устанавливаться одним из следующих способов: 
\begin{itemize}
\item на основании внешнего сигнала от сборки блок-контактов выключателя;
\item функцией фиксации положения выключателя (ФПВ), которая пофазно контролирует состояние блок-контактов выключателя и в случае расхождения полюсов формирует сигнал \textit{«Внутр. пуск ЗНФ»}.
\end{itemize}

\item
В зависимости от используемого способа соответствующий сигнал должен назначаться на конфигурируемый вход \textit{«Неперекл. фаз»}.  

\item
Действие на отключение выключателя с запретом АПВ осуществляется с выдержками времени \textbf{«Тср ЗНФ»}. Задаваемые выдержки времени предназначена для отстройки от разновременности переключения блок-контактов выключателя.

\item
Предусмотрена возможность блокировки ЗНФ в цикле ОАПВ \textit{«ЗНФ: Цикл ОАПВ»} на время, которое задаётся уставкой \textbf{«Тпрод бл.ЗНФ от ОАПВ»}.

\item
Оперативный ввод/вывод ЗНФ производится с помощью БУ (блока управления функцией) \textit{«Вывод ЗНФ»}.

\AbzazOne

\vspace{2mm}
\begin{figure}[H]
\centering
\includegraphics[width=0.9\textwidth,height=0.9\textheight,keepaspectratio]{\fbpath/11.01 ШР/761.3441 ЗНФ/img/ZNF.pdf}
\captionof{figure}{Функциональная схема логики ЗНФ}\label{pict:ZNF}
\end{figure}


\end{enumerate} 